\documentclass[12pt]{article}
\usepackage[T1]{fontenc}
\usepackage[utf8]{inputenc}
\usepackage{lmodern}
\usepackage{textcomp}
\usepackage{enumitem}
\usepackage{geometry}
\geometry{margin=1in}
\begin{document}
\begin{center}
Answer Key - History - Undergraduate Level Assessment 25 \
\small (Answers are ordered exactly as in the question paper.)
\end{center}

\section*{Multiple Choice}
\begin{enumerate}[label=\arabic*.]
\item \textbf{Answer:} A
\\ \textit{Options:} A) stratigraphic analysis; B) resistivity survey; C) metal detectors; D) ground-penetrating radar; E) underwater archaeology
\\ \textit{Note:} Stratigraphic analysis is the most effective archaeological method for identifying large-scale land modifications because it involves examining the layers of soil and artifacts deposited over time, allowing researchers to understand the sequence and extent of human alterations to the landscape.
\item \textbf{Answer:} A
\\ \textit{Options:} A) Caral; B) Chichen Itza; C) The Parthenon; D) Machu Picchu; E) Angkor Wat
\\ \textit{Note:} Caral, an ancient city in present-day Peru, is known for its monumental earthworks and distinctive semicircular concentric ridges that date back to approximately 3,500 years ago, making it a significant archaeological site in the study of early urban development and culture in the Americas.
\item \textbf{Answer:} A
\\ \textit{Options:} A) Egypt; B) Western Europe; C) North America; D) South America; E) Eastern China
\\ \textit{Note:} The world's first civilization developed in Egypt, particularly along the fertile banks of the Nile River, where the combination of agricultural innovation, stable food supply, and social organization led to the rise of complex societies around 3100 BCE.
\item \textbf{Answer:} A
\\ \textit{Options:} A) pelagic; lacustrine; B) pelagic; littoral; C) pelagic; midden; D) littoral; lacustrine; E) midden; littoral
\\ \textit{Note:} The correct answer is based on the distinction between the Lake Forest Archaic tradition, which primarily utilized pelagic resources from open waters, and the Maritime Archaic, which focused on hunting lacustrine creatures found in lakes and freshwater environments.
\item \textbf{Answer:} C
\\ \textit{Options:} A) the development of language and writing.; B) slowly increased brain size and the first stone tools.; C) a rapid increase in brain size.; D) the transition from bipedalism to quadrupedalism.; E) a decrease in stature and increase in brain size.
\\ \textit{Note:} After 400,000 years ago, hominid evolution is characterized by a rapid increase in brain size, which is supported by archaeological evidence indicating that significant advancements in tool-making, social structures, and cognitive abilities emerged alongside this increase.
\end{enumerate}

\section*{True / False}
\begin{enumerate}[label=\arabic*.]
\setcounter{enumi}{5}
\item \textbf{Answer:} True
\\ \textit{Options:} A) True; B) False
\\ \textit{Note:} According to the passage: 1821
\item \textbf{Answer:} False
\\ \textit{Options:} A) True; B) False
\\ \textit{Note:} The correct answer is: speak well and write
\item \textbf{Answer:} False
\\ \textit{Options:} A) True; B) False
\\ \textit{Note:} The correct answer is: 10,550,350
\item \textbf{Answer:} True
\\ \textit{Options:} A) True; B) False
\\ \textit{Note:} According to the passage: 2006
\item \textbf{Answer:} True
\\ \textit{Options:} A) True; B) False
\\ \textit{Note:} According to the passage: Albert Grzymała
\end{enumerate}

\section*{Long Form Response}
\begin{enumerate}[label=\arabic*.]
\setcounter{enumi}{10}
\item \textbf{Answer:} industry dispersal may have been a factor
\\ \textit{Note:} Expected answer: industry dispersal may have been a factor
\item \textbf{Answer:} 1951
\\ \textit{Note:} Expected answer: 1951
\end{enumerate}

\vfill
\noindent\textit{End of Answer Key}
\end{document}